%% ============================================================
%%  handout-aula11.tex — Resumo de estudo da Aula 11
%%  Bandits bayesianos e amostragem de Thompson
%%  Aprendizagem por Reforço · UFU · Prof. Marcelo Keese Albertini
%%  Adaptado de CS234 (Stanford), lecture11post.pdf
%%  Compilar com XeLaTeX (2x) — precisa de handoutAprendizadoRL.sty
%% ============================================================
\documentclass[11pt, a4paper]{article}
%% .sty do projeto na raiz do repositório (../)
\makeatletter
\def\input@path{{./}{../}}
\makeatother


\usepackage{handoutAprendizadoRL}   % antes do babel (carrega fontspec)
\usepackage{babel}
\babelprovide[import, main]{brazilian}
\usepackage{multicol}

\setaula{Aula 11 — Bandits bayesianos}
\setprof{Prof.\ Marcelo Keese Albertini}

%% ---- macros locais
\newcommand{\arrep}{\mathrm{Arrep}}
\newcommand{\arrepb}{\mathrm{ArrepBayes}}
\newcommand{\alg}{\mathcal{A}}
\newcommand{\Bdist}{\operatorname{Beta}}
\newcommand{\Bern}{\operatorname{Bernoulli}}
\newcommand{\Nrm}{\mathcal{N}}
\newcommand{\hist}{h_t}
\newcommand{\Ind}{\mathbb{1}}
\newcommand{\subt}[1]{\medskip\noindent\textbf{\color{rlAzulClaro}#1}\par\nobreak\smallskip}

\hyphenation{apren-di-za-gem re-com-pen-sa se-quen-cial de-ter-mi-nís-ti-co
             ar-re-pen-di-men-to ba-ye-si-a-no con-ju-ga-ção pos-te-ri-or
             uni-di-men-si-o-nais uni-di-men-si-o-nal in-de-pen-den-tes
             im-plau-sí-veis con-tex-tu-al con-tex-tu-ais pseu-do-con-ta-gens
             hi-per-pa-râ-me-tro con-cen-tra-ção pro-ba-bi-li-da-de}

\begin{document}

\capa{Aula 11 — Bandits bayesianos e amostragem de Thompson}
     {Resumo de estudo \ \textbullet\ \ conjugação Beta--Bernoulli, casamento de
      probabilidade e arrependimento bayesiano}
     {Prof.\ Marcelo Keese Albertini}
     {Adaptado e expandido de CS234: Reinforcement Learning, Emma Brunskill (Stanford)}

\begin{ideiabox}
A aula anterior resolveu a exploração por \emph{otimismo}: aja como se o melhor
mundo compatível com seus dados fosse o verdadeiro. Esta aula apresenta a
segunda grande família: aja como se \emph{um mundo sorteado ao acaso} entre os
compatíveis fosse o verdadeiro. A segunda ideia cabe em seis linhas de código,
não tem hiperparâmetro de exploração, e tem essencialmente as mesmas garantias
teóricas.
\end{ideiabox}

%% ============================================================
\section{O problema, em uma página}

\begin{defbox}{bandit multiarmado}
Par $(\gA, \mathcal{R})$ com $\gA$ um conjunto \emph{conhecido} de $m$ ações
(braços) e $\mathcal{R}^a(r) = \Prob[r \mid a]$ uma distribuição
\emph{desconhecida} de recompensas para cada braço. A cada passo $t$ o agente
escolhe $a_t \in \gA$ e recebe $r_t \sim \mathcal{R}^{a_t}$. Escreve-se
$Q(a) = \E[r \mid a]$ e $\Vstar = \max_a Q(a)$.
\end{defbox}

Um bandit é um MDP de um único estado. A ação não move o agente; ela só produz
recompensa e informação. Com isso, a equação de Bellman degenera em
$Q(a) = \E[r\mid a]$ e todo o conteúdo do problema passa a ser \emph{exploração}
— por isso o bandit é o laboratório onde as ideias de exploração são inventadas
antes de subirem para o MDP completo.

\begin{defbox}{lacuna, arrependimento por passo e arrependimento total}
$\Delta_a = \Vstar - Q(a)$ é a \emph{lacuna} do braço $a$;
$\ell_t = \E[\Vstar - Q(a_t)]$ é o arrependimento do passo $t$; e
\[
  L_T \;=\; \E\!\left[\sum_{t=1}^{T} \Vstar - Q(a_t)\right]
\]
é o arrependimento total. Maximizar recompensa acumulada equivale a minimizar
$L_T$. O objetivo é \emph{sublinearidade}: $L_T/T \to 0$.
\end{defbox}

\subt{A identidade que organiza toda a análise}

Seja $N_T(a) = \sum_{t=1}^{T}\Ind(a_t = a)$ o número de puxadas do braço $a$.
Trocando a ordem da soma:
\begin{equation}
  L_T \;=\; \E\!\left[\sum_{t=1}^{T}\sum_{a \in \gA} \Ind(a_t = a)\,\Delta_a\right]
      \;=\; \sum_{a \in \gA} \E[N_T(a)]\;\Delta_a .
  \label{eq:decomp}
\end{equation}

\begin{ideiabox}
A equação~\eqref{eq:decomp} diz que \textbf{arrependimento é contagem de erros
ponderada pela gravidade de cada erro}. Ela transforma toda demonstração de cota
de arrependimento em um problema de contagem: basta mostrar que
$\E[N_T(a)]$ é pequeno para os braços com $\Delta_a$ grande. Guarde-a — ela é
usada na demonstração do UCB, na de Thompson e na cota inferior de Lai e
Robbins.
\end{ideiabox}

\subt{Três limitações do otimismo que motivam esta aula}

\begin{enumerate}
  \item \textbf{UCB é determinístico.} Dado o histórico, a ação está decidida.
        Se o retorno chega em lote — mil visitantes num site antes de sabermos
        se o primeiro clicou — o índice não muda e o algoritmo repete a mesma
        escolha mil vezes.
  \item \textbf{UCB ignora conhecimento prévio.} Se dez anos de prontuários
        indicam que um tratamento cura cerca de $90\%$ dos casos, o UCB começa
        do zero mesmo assim.
  \item \textbf{O bônus $\sqrt{2\ln t / N_t(a)}$ é genérico.} Ele vale para
        qualquer distribuição limitada, e por isso é conservador demais para a
        distribuição que você de fato tem.
\end{enumerate}

\begin{discbox}
Num bandit com recompensas \emph{determinísticas} — cada braço devolve sempre o
mesmo valor —, o que acontece com o UCB? Ele tem arrependimento sublinear com
probabilidade $1$, ou existe probabilidade estritamente positiva de
arrependimento linear? E a divisão por $N_t(a)$ chega a ser um problema?

\smallskip
{\footnotesize\color{rlCinza}\emph{Direção da resposta:} examinar por que o termo
de confiança existe — cobrir o ruído amostral da média empírica — e o que
sobra desse papel quando a variância é zero. Vale também discutir a
inicialização, em que cada braço é puxado uma vez antes de o índice ser
avaliado.}
\end{discbox}

%% ============================================================
\section{Inferência bayesiana: só o necessário}

Na visão bayesiana, partimos de um \emph{prior} sobre os parâmetros
desconhecidos — aqui, sobre a distribuição de recompensas de cada braço — e o
atualizamos com os dados pela regra de Bayes:
\begin{equation}
  \underbrace{p(\theta \mid r)}_{\text{posterior}}
  \;=\;
  \frac{\overbrace{p(r \mid \theta)}^{\text{verossimilhança}}\;
        \overbrace{p(\theta)}^{\text{prior}}}
       {\underbrace{\int p(r\mid\theta)p(\theta)\,d\theta}_{\text{evidência}}}
  \;\;\propto\;\; p(r \mid \theta)\,p(\theta).
  \label{eq:bayes}
\end{equation}

A mudança de objeto é o ponto: em vez da estimativa pontual $\hat{Q}(a)$ com uma
barra de erro colada por fora, carregamos uma \emph{distribuição inteira} sobre
$Q(a)$. A incerteza deixa de ser um apêndice do algoritmo e passa a ser o estado
de conhecimento do agente.

\begin{alertabox}
O denominador de~\eqref{eq:bayes} é uma integral que, sem estrutura adicional,
não tem forma fechada — restaria MCMC ou inferência variacional. Como o bandit
faz \emph{uma} atualização por passo, isso seria inviável. Conjugação não é
elegância matemática: é requisito de engenharia.
\end{alertabox}

\begin{defbox}{prior conjugado}
$p(\theta)$ é \emph{conjugado} para a verossimilhança $p(r\mid\theta)$ se o
posterior $p(\theta \mid r)$ pertence à mesma família paramétrica do prior. A
atualização vira então um mapa entre hiperparâmetros — tipicamente somas — com
custo e memória $O(1)$ por observação. Toda família exponencial admite prior
conjugado.
\end{defbox}

\begin{center}\small
\begin{tabular}{@{}lll@{}}
\toprule
\textbf{Recompensa} & \textbf{Prior conjugado} & \textbf{Uso típico} \\
\midrule
$\Bern(\theta)$                    & $\Bdist(\alpha,\beta)$ & clique, cura/falha, conversão \\
$\Nrm(\mu,\sigma^2)$, $\sigma$ conhecido & $\Nrm(\mu_0,\sigma_0^2)$ & recompensa contínua \\
Categórica em $K$ valores          & $\mathrm{Dirichlet}(\alpha_1,\ldots,\alpha_K)$ & nota de $1$ a $5$ \\
$\mathrm{Poisson}(\lambda)$        & $\mathrm{Gama}(\alpha,\beta)$ & contagens \\
\bottomrule
\end{tabular}
\end{center}

\subsection*{O par Beta--Bernoulli}

\begin{defbox}{distribuição Beta}
\[
  p(\theta \mid \alpha,\beta)
  = \frac{\Gamma(\alpha+\beta)}{\Gamma(\alpha)\Gamma(\beta)}\,
    \theta^{\alpha-1}(1-\theta)^{\beta-1},
  \qquad \theta \in [0,1],\ \ \alpha,\beta>0,
\]
com $\E[\theta] = \frac{\alpha}{\alpha+\beta}$, moda
$\frac{\alpha-1}{\alpha+\beta-2}$ (para $\alpha,\beta>1$) e
$\Var[\theta] = \frac{\alpha\beta}{(\alpha+\beta)^2(\alpha+\beta+1)}$.
Em particular $\Bdist(1,1)$ é a uniforme em $[0,1]$.
\end{defbox}

\begin{teobox}{Proposição — conjugação Beta--Bernoulli}
Se $\theta \sim \Bdist(\alpha,\beta)$ e observamos $r \in \{0,1\}$ com
$r \sim \Bern(\theta)$, então $\theta \mid r \sim \Bdist(\alpha+r,\ \beta+1-r)$.
\end{teobox}

\begin{provabox}
Pela proporcionalidade em~\eqref{eq:bayes}, basta acompanhar os expoentes:
\[
  p(\theta\mid r)
  \;\propto\;
  \underbrace{\theta^{r}(1-\theta)^{1-r}}_{\text{verossimilhança }\Bern}
  \cdot
  \underbrace{\theta^{\alpha-1}(1-\theta)^{\beta-1}}_{\text{prior }\Bdist}
  \;=\;
  \theta^{(\alpha+r)-1}(1-\theta)^{(\beta+1-r)-1}.
\]
O lado direito é o núcleo de uma $\Bdist(\alpha+r,\ \beta+1-r)$. Como a família
Beta é normalizável e o posterior é uma densidade, a constante de normalização
está determinada e não precisa ser calculada. $\square$

\smallskip
Por indução, após $S$ sucessos e $F$ fracassos em qualquer ordem, o posterior é
$\Bdist(\alpha+S,\ \beta+F)$ — a ordem das observações é irrelevante, o que era
de se esperar, já que as recompensas são i.i.d.\ dado $\theta$.
\end{provabox}

\begin{ideiabox}
\textbf{Leitura em pseudo-contagens.} $\alpha-1$ e $\beta-1$ são sucessos e
fracassos \emph{imaginários} que o prior injeta antes de qualquer dado.
Partindo de $\Bdist(1,1)$, a média posterior após $S$ sucessos e $F$ fracassos é
\[
  \E[\theta \mid S,F] \;=\; \frac{S+1}{S+F+2},
\]
a chamada regra de sucessão de Laplace: dois dados fantasmas que impedem o
agente de concluir ``$100\%$ de acerto'' depois de um único acerto. A variância
cai como $1/(\alpha+\beta)$: mais dados, posterior mais estreito.
\end{ideiabox}

\begin{center}
\begin{tikzpicture}[x=6.4cm, y=0.30cm]
  \draw[rlCinza!55] (0,0) -- (1.03,0);
  \foreach \x in {0,0.25,0.5,0.75,1}{
    \draw[rlCinza!55] (\x,0) -- ++(0,-2pt)
      node[below, font=\tiny, color=rlCinza, inner sep=1pt] {\x};}
  \node[font=\scriptsize, color=rlCinza, anchor=west] at (1.05,0) {$\theta$};
  \draw[rlCinza!45, thick] plot[smooth] coordinates
    {(0.000,1.000) (0.500,1.000) (1.000,1.000)};
  \draw[rlAzul!45, thick] plot[smooth] coordinates
    {(0.000,0.000) (0.114,0.006) (0.227,0.047) (0.341,0.158) (0.455,0.376)
     (0.568,0.734) (0.682,1.268) (0.795,2.013) (0.886,2.785) (0.955,3.479) (1.000,4.000)};
  \draw[rlTeal, thick] plot[smooth] coordinates
    {(0.000,0.000) (0.205,0.001) (0.318,0.016) (0.409,0.082) (0.500,0.281)
     (0.591,0.741) (0.682,1.569) (0.750,2.403) (0.818,3.213) (0.864,3.518)
     (0.909,3.359) (0.955,2.363) (0.977,1.393) (1.000,0.000)};
  \draw[rlLaranja, very thick] plot[smooth] coordinates
    {(0.000,0.000) (0.500,0.000) (0.591,0.003) (0.659,0.036) (0.705,0.164)
     (0.750,0.609) (0.795,1.837) (0.841,4.332) (0.886,7.266) (0.909,7.752)
     (0.932,6.693) (0.955,3.989) (0.977,0.986) (1.000,0.000)};
  \node[font=\tiny, color=rlCinza,   anchor=south] at (0.12,1.05) {$\Bdist(1,1)$};
  \node[font=\tiny, color=rlAzul!55, anchor=south east] at (0.99,3.9) {$\Bdist(4,1)$};
  \node[font=\tiny, color=rlTeal,    anchor=south east] at (0.79,3.3) {$\Bdist(8,2)$};
  \node[font=\tiny, color=rlLaranja, anchor=south east] at (0.90,7.7) {$\Bdist(30,4)$};
\end{tikzpicture}
\end{center}

\noindent
{\footnotesize\color{rlCinza} O posterior de um braço ao longo de trinta e três
observações. A massa migra para a direita \emph{e} encolhe: a média sobe de
$0{,}5$ para $30/34 \approx 0{,}88$ e o desvio-padrão cai de $0{,}29$ para
$0{,}055$. A largura da curva \emph{é} a incerteza — e é ela que dosará a
exploração.}

%% ============================================================
\section{Bandits bayesianos}

\begin{defbox}{bandit bayesiano}
Bandit em que se assume conhecimento prévio $p[\mathcal{R}]$ sobre as
recompensas. A cada passo computa-se o posterior $p[\mathcal{R}\mid \hist]$ com
$\hist = (a_1,r_1,\ldots,a_{t-1},r_{t-1})$, e o posterior guia a exploração.
No caso Bernoulli, todo o estado de conhecimento cabe em $2m$ números
$(\alpha_a, \beta_a)$.
\end{defbox}

Há duas maneiras clássicas de transformar um posterior em uma ação.

\subt{Bayes-UCB: otimismo calibrado}
Escolha um quantil alto do posterior,
$a_t = \argmax_a q_{1-1/t}\big(p(Q(a)\mid \hist)\big)$. A ideia é a mesma do UCB,
mas o ``bônus'' vem da forma real do posterior, e não de uma desigualdade de
concentração válida para o pior caso. (Kaufmann, Cappé e Garivier, 2012.)

\subt{Casamento de probabilidade}
Escolha cada ação com a probabilidade de que ela seja a ótima:
\begin{equation}
  \pol(a \mid \hist) \;=\; \Prob\big[\,Q(a) > Q(a'),\ \forall a'\neq a \ \big|\ \hist\,\big].
  \label{eq:pm}
\end{equation}
Ações incertas têm caudas mais gordas e, portanto, maior probabilidade de serem
o máximo: a exploração emerge da definição, sem bônus algum. O problema é
computar~\eqref{eq:pm}, uma integral $m$-dimensional sem forma fechada simples.

\begin{alertabox}
\textbf{Cuidado com o nome.} Em psicologia experimental,
``\emph{probability matching}'' designa escolher a alternativa $A$ numa fração
das vezes igual à \emph{frequência com que $A$ é recompensada} — comportamento
comprovadamente subótimo. Aqui a probabilidade é outra: a de que $A$ seja o
\emph{melhor braço}, dada a evidência acumulada. Os dois procedimentos só
coincidem no caso degenerado de um braço único.
\end{alertabox}

%% ============================================================
\section{Amostragem de Thompson}

\begin{algbox}{amostragem de Thompson (Thompson, 1933)}
\begin{enumerate}[leftmargin=1.6em, itemsep=0.15em]
  \item Inicialize um prior $p(\mathcal{R}^a)$ para cada braço $a$.
  \item Para $t = 1, 2, \ldots$:
  \begin{enumerate}[leftmargin=1.4em, label=(\alph*), itemsep=0.1em]
    \item para cada braço $a$, \emph{sorteie}
          $\tilde{\mathcal{R}}^a \sim p(\mathcal{R}^a \mid \hist)$ e calcule
          $\tilde Q(a) = \E[\tilde{\mathcal{R}}^a]$;
    \item jogue $a_t = \argmax_{a\in\gA} \tilde Q(a)$;
    \item observe $r_t$ e atualize o posterior de $a_t$ pela regra de Bayes.
  \end{enumerate}
\end{enumerate}
No caso Bernoulli com prior Beta, o passo (a) é
$\tilde\theta_a \sim \Bdist(\alpha_a,\beta_a)$ e o passo (c) soma $1$ em
$\alpha_{a_t}$ (se $r_t=1$) ou em $\beta_{a_t}$ (se $r_t=0$).
\end{algbox}

\begin{ideiabox}
Leia o passo (a) como: \emph{``finja, por um instante, que um dos mundos
compatíveis com meus dados é o mundo verdadeiro — e aja de forma ótima nele''}.
Daí o nome alternativo, \textbf{amostragem posterior}.
\end{ideiabox}

\begin{teobox}{Proposição — Thompson implementa casamento de probabilidade}
A probabilidade de a amostragem de Thompson jogar o braço $a$ no passo $t$ é
exatamente $\pol(a\mid \hist)$, definida em~\eqref{eq:pm}.
\end{teobox}

\begin{provabox}
O braço jogado é $\argmax_a \tilde Q(a)$, com $\tilde Q$ sorteado do posterior.
Logo
\[
  \Prob[a_t = a]
  \;=\;
  \E_{\tilde Q \sim p(\cdot\mid \hist)}
    \Big[\Ind\big(a = \argmax\nolimits_{a'\in\gA}\tilde Q(a')\big)\Big]
  \;=\;
  \Prob\big[Q(a) > Q(a'),\ \forall a'\neq a \mid \hist\big]
  \;=\;
  \pol(a\mid \hist),
\]
onde a segunda igualdade é apenas a definição de esperança de uma função
indicadora sob o posterior. $\square$
\end{provabox}

\begin{ideiabox}
O conteúdo da proposição, em uma frase: \textbf{sortear do posterior e tomar o
$\argmax$ é um estimador de Monte Carlo com uma \emph{única} amostra} da
integral que não sabíamos calcular. E uma amostra basta, porque não precisamos
do \emph{valor} de $\pol(a\mid\hist)$ — precisamos apenas de uma ação
\emph{distribuída} segundo ele. É uma das trocas mais elegantes da área: o
problema difícil não foi resolvido, foi contornado.
\end{ideiabox}

\subsection*{Um exemplo numérico completo}

Três tratamentos para um dedo quebrado, com parâmetros verdadeiros
(desconhecidos do agente) $\theta_1 = 0{,}95$ para a cirurgia,
$\theta_2 = 0{,}90$ para o esparadrapo e $\theta_3 = 0{,}10$ para não fazer
nada. Prior $\Bdist(1,1)$ nos três braços.

\begin{center}\small
\begin{tabular}{@{}clccl@{}}
\toprule
\textbf{Passo} & \textbf{Posteriores antes} & \textbf{Amostras} & \textbf{Escolha} & \textbf{Resultado e atualização} \\
\midrule
1 & $(1,1)\;(1,1)\;(1,1)$ & $0{,}30\ \ 0{,}50\ \ 0{,}60$ & $a_3$ & $r=0 \Rightarrow \theta_3 \sim \Bdist(1,2)$ \\
2 & $(1,1)\;(1,1)\;(1,2)$ & $0{,}70\ \ 0{,}50\ \ 0{,}30$ & $a_1$ & $r=1 \Rightarrow \theta_1 \sim \Bdist(2,1)$ \\
3 & $(2,1)\;(1,1)\;(1,2)$ & $0{,}71\ \ 0{,}65\ \ 0{,}10$ & $a_1$ & $r=1 \Rightarrow \theta_1 \sim \Bdist(3,1)$ \\
4 & $(3,1)\;(1,1)\;(1,2)$ & $0{,}75\ \ 0{,}45\ \ 0{,}40$ & $a_1$ & $r=1 \Rightarrow \theta_1 \sim \Bdist(4,1)$ \\
\bottomrule
\end{tabular}
\end{center}

\begin{alertabox}
Exemplo \emph{fictício}: estes não são os desfechos reais das opções de
tratamento para um dedo quebrado. O caso serve apenas para acompanhar a
aritmética.
\end{alertabox}

Na mesma instância, o otimismo puxou a sequência
$a_1$, $a_2$, $a_3$, $a_1$, $a_2$, e Thompson puxou
$a_3$, $a_1$, $a_1$, $a_1$, $a_1$. A diferença é qualitativa: o otimismo
\emph{tem de} experimentar cada braço até que sua incerteza caia, porque o
bônus é enorme para quem tem poucas puxadas; Thompson descarta braços que a
evidência já tornou implausíveis, mesmo com um único dado. Numa amostra o TS
ganha, noutra perde — o que se compara é a média.

\subsection*{Implementação}

\begin{algbox}{Thompson Bernoulli, em seis linhas}
\begin{verbatim}
alfa = np.ones(K); beta = np.ones(K)   # prior Beta(1,1) por braco
for t in range(T):
    theta = rng.beta(alfa, beta)       # 1 amostra de cada posterior
    a = int(np.argmax(theta))          # otimo no mundo sorteado
    r = puxar(a)                       # recompensa em {0,1}
    alfa[a] += r; beta[a] += 1 - r     # atualizacao conjugada
\end{verbatim}
\end{algbox}

Nenhum hiperparâmetro de exploração, nenhum cronograma, nenhuma cota de
concentração. Custo por passo: $O(m)$ em tempo e memória — o mesmo do
$\epsilon$-guloso. Para recompensas contínuas, trocam-se as duas últimas linhas
pela atualização Normal--Normal; a estrutura do laço não muda.

\subsection*{O que acontece na média}

Simulação com $\theta = (0{,}95;\ 0{,}90;\ 0{,}10)$ e prior $\Bdist(1,1)$,
com média sobre $1500$ execuções independentes:

\begin{center}\small
\begin{tabular}{@{}lrrrr@{}}
\toprule
\textbf{Arrependimento acumulado em} & $T=100$ & $T=500$ & $T=1000$ & $T=2000$ \\
\midrule
Amostragem de Thompson   & $3{,}2$ & $6{,}4$  & $7{,}9$  & $9{,}2$ \\
Guloso                   & $2{,}6$ & $7{,}8$  & $14{,}1$ & $26{,}5$ \\
UCB                      & $7{,}3$ & $19{,}5$ & $29{,}9$ & $45{,}7$ \\
$\epsilon$-guloso ($\epsilon = 0{,}1$) & $5{,}1$ & $20{,}0$ & $36{,}9$ & $68{,}6$ \\
\bottomrule
\end{tabular}
\end{center}

\begin{alertabox}
\textbf{Não conclua que o guloso é bom.} O guloso tem arrependimento
\emph{linear}: em $23\%$ das execuções ele trava no braço $a_2$ e nunca mais
sai. Ele só parece competitivo porque $\Delta_2 = 0{,}05$ é minúsculo — a reta
tem inclinação pequena, mas é reta. Estendendo o horizonte, seu arrependimento
médio vai a $13{,}2$ ($T=10^3$), $58{,}9$ ($T=5\times10^3$) e $229{,}9$
($T=2\times10^4$), enquanto o de Thompson vai a $9{,}2$, $13{,}0$ e $15{,}7$ —
crescimento logarítmico. \emph{Moral metodológica:} nunca compare algoritmos de
exploração em um único horizonte, nem olhando apenas a média.
\end{alertabox}

%% ============================================================
\section{Arrependimento bayesiano}

\begin{defbox}{arrependimento frequentista e bayesiano}
\emph{Frequentista} — supõe parâmetros verdadeiros e fixos $\theta$, com a
esperança tomada sobre o histórico gerado pelo algoritmo $\alg$:
\[
  \arrep(\alg,T;\theta) \;=\; \E_{\alg}\!\left[\sum_{t=1}^{T} Q(a^{*}) - Q(a_t)\ \Big|\ \theta\right].
\]
\emph{Bayesiano} — supõe um prior $p_\theta$ e tira a média também sobre ele:
\[
  \arrepb(\alg,T) \;=\; \E_{\theta \sim p_\theta}\big[\arrep(\alg,T;\theta)\big].
\]
\end{defbox}

\begin{teobox}{Observação — uma cota implica a outra}
$\arrepb(\alg,T) \;\le\; \sup_{\theta}\,\arrep(\alg,T;\theta)$, para todo prior
$p_\theta$.
\end{teobox}

\begin{provabox}
A média de uma função é no máximo o seu supremo:
$\E_{\theta\sim p_\theta}[\arrep(\alg,T;\theta)] \le \sup_\theta \arrep(\alg,T;\theta)$.
$\square$

\smallskip
Consequência prática: \textbf{toda cota frequentista uniforme já é uma cota
bayesiana}. A recíproca é falsa — uma cota bayesiana boa pode conviver com
desempenho péssimo em um $\theta$ de probabilidade a priori desprezível. É
exatamente aí que mora a diferença de dificuldade entre as duas análises, e é
por isso que uma cota \emph{frequentista} para Thompson é um resultado bem mais
forte (e bem mais recente) do que a bayesiana.
\end{provabox}

\subt{Otimismo como técnica de demonstração}

Suponha que $U_t(a)$ seja, com alta probabilidade, uma cota superior de $Q(a)$.
No evento em que a cota vale,
\begin{equation}
  \arrep(\alg,T;\theta)
  = \E\!\left[\sum_{t=1}^{T} Q(a^{*}) - Q(a_t)\right]
  \;\le\;
  \E\!\left[\sum_{t=1}^{T} \mdestaque{rlLaranja}{U_t(a_t) - Q(a_t)}\right].
  \label{eq:otim}
\end{equation}
O passo usa duas coisas: $Q(a^*) \le U_t(a^*)$, porque a cota vale para o braço
ótimo; e $U_t(a^*) \le U_t(a_t)$, porque o algoritmo escolheu $a_t$ justamente
por ter o maior índice. À direita não sobra nenhuma referência a $a^*$: o
arrependimento virou a soma das larguras dos intervalos de confiança dos braços
\emph{efetivamente puxados} — e cotar isso é contar puxadas, via~\eqref{eq:decomp}.

\begin{ideiabox}
Eis a razão profunda de o otimismo ser onipresente na teoria de bandits: ele não
é apenas uma heurística de exploração, é o artifício que faz a demonstração
fechar. O notável é que a análise bayesiana de Thompson usa \emph{esse mesmo
artifício} — via a chamada \emph{razão de informação} de Russo e Van Roy —
apesar de o algoritmo não ser otimista. Thompson herda as cotas do otimismo sem
ser otimista.
\end{ideiabox}

\begin{center}\small
\begin{tabular}{@{}llp{6.4cm}@{}}
\toprule
\textbf{Garantia} & \textbf{Cota} & \textbf{Referência e observação} \\
\midrule
UCB, frequentista & $O(\sqrt{mT\ln T})$ &
  Auer, Cesa-Bianchi e Fischer (2002); também
  $O\big(\sum_a \ln T/\Delta_a\big)$, dependente das lacunas \\
Thompson, bayesiano & $O(\sqrt{mT\ln T})$ &
  Russo e Van Roy (2014), via razão de informação; vale para qualquer prior \\
Thompson, frequentista & $O(\sqrt{mT\ln T})$ &
  Agrawal e Goyal (2012, 2013), Bernoulli com prior Beta; constantes piores
  que as do UCB \\
Cota inferior (qualquer algoritmo) & $\Omega(\sqrt{mT})$ &
  Lai e Robbins (1985); Auer et al.\ (2002) \\
\bottomrule
\end{tabular}
\end{center}

\medskip
\noindent
Em ordem de grandeza, portanto, \textbf{empate}: amostragem posterior tem as
mesmas cotas do UCB a menos de constantes, e ambos são ótimos em ordem. Nas
constantes e na prática o TS costuma vencer, sobretudo em bandits contextuais,
onde calibrar o bônus de otimismo é difícil.

%% ============================================================
\section{Onde Thompson falha, e o que vem depois}

\subt{Prior enganoso}
As garantias bayesianas são \emph{em média sobre o prior}. Se o prior estiver
errado, não há promessa alguma. Com dois braços ($\theta_{\text{ruim}} = 0{,}10$
e $\theta_{\text{bom}} = 0{,}50$) e $T = 1000$:

\begin{center}\small
\begin{tabular}{@{}lcc@{}}
\toprule
\textbf{Prior do braço ruim} & \textbf{Puxadas do braço ruim} & \textbf{Arrependimento} \\
\midrule
$\Bdist(1,1)$ — uniforme   & $13{,}6$  & $5{,}4$ \\
$\Bdist(100,1)$ — enganoso & $204{,}0$ & $81{,}6$ \\
\bottomrule
\end{tabular}
\end{center}

\noindent
$\Bdist(100,1)$ afirma, antes de qualquer dado, que o braço acerta cerca de
$99\%$ das vezes. Trazê-lo de volta à realidade exige acumular cerca de cem
fracassos, e cada um deles é um passo desperdiçado: o arrependimento fica quinze
vezes maior. \emph{Prática recomendada:} priors fracos ($\alpha+\beta$ pequeno),
salvo evidência histórica real. Um prior informativo é uma afirmação empírica e
deve ser defendida como tal.

\subt{\emph{Feedback} em lote}
Quando o retorno demora, o UCB — sendo determinístico — repete a mesma ação até
que o histórico mude. Thompson, sorteando a cada decisão, \emph{diversifica
naturalmente} dentro do lote: mil decisões tomadas com o mesmo posterior
produzem uma mistura de braços na proporção $\pol(a\mid \hist)$. Foi um dos
motivos empíricos do sucesso do TS em recomendação em larga escala.

\subt{Bandits contextuais}
No bandit contextual, o ambiente sorteia i.i.d.\ um contexto $x_t$ que afeta a
recompensa de cada braço; o agente vê $x_t$, escolhe $a_t$ e recebe
$r_t \sim \mathcal{R}(\cdot \mid x_t, a_t)$. É o passo intermediário entre
bandit e MDP: já há generalização entre situações, mas a ação ainda não move o
estado. O TS se estende sem esforço — mantenha um posterior sobre os
\emph{parâmetros} de $Q_w(x,a)$, sorteie $\tilde w$, jogue
$\argmax_a Q_{\tilde w}(x_t,a)$. Chapelle e Li (2010) mostraram que, em
recomendação de notícias, o TS iguala ou supera o LinUCB, com implementação mais
simples e maior robustez a atrasos.

\subt{A política ótima e o índice de Gittins}
Dado um prior e um horizonte, existe uma política que maximiza a recompensa
esperada, e TS não é ela. Formalmente, o problema bayesiano é um MDP cujo estado
é a \emph{crença} — todos os pares $(\alpha_a,\beta_a)$ —, cuja ação é o braço e
cuja transição é a atualização de Bayes. Chama-se \emph{MDP de crenças}, e a
política ótima sai de programação dinâmica sobre ele. Computacionalmente isso é
inviável de forma ingênua: o conjunto de crenças alcançáveis cresce
combinatoriamente com o horizonte.

\begin{defbox}{política de índice}
Política que calcula um índice de valor real para cada braço usando somente as
estatísticas \emph{daquele braço} e o horizonte, e joga o braço de maior índice.
(Definição de Lattimore e Szepesvári, \emph{Bandit Algorithms}.)
\end{defbox}

\begin{teobox}{Teorema do índice (Gittins, 1979)}
Para um bandit bayesiano com braços \emph{independentes} e critério de recompensa
esperada \emph{descontada} em horizonte infinito, a política ótima é uma
política de índice: jogue sempre o braço de maior índice de Gittins.
\end{teobox}

O resultado é notável: reduz um problema de dimensão $m$ a $m$ problemas
unidimensionais independentes. As limitações explicam por que ele não é usado
na prática — exige braços independentes, quebra com horizonte finito ou com o
critério de recompensa média, o índice é caro de calcular, e não há extensão
para o caso contextual.

\begin{ideiabox}
Note a inversão conceitual do parágrafo anterior. O bandit começou como o caso
\emph{fácil} do MDP — um único estado. Ao exigirmos otimalidade sob incerteza,
ele voltou a ser um MDP, agora de estado contínuo e alta dimensão.
\textbf{A incerteza é, ela própria, um estado.}
\end{ideiabox}

\subt{De volta ao RL}
A ideia de Thompson não depende de o problema ser um bandit. O padrão geral é:
\emph{sorteie um mundo plausível, aja de forma ótima nele por um episódio
inteiro, atualize}.

\begin{center}\small
\begin{tabular}{@{}lp{7.0cm}l@{}}
\toprule
\textbf{Método} & \textbf{O que se sorteia} & \textbf{Referência} \\
\midrule
PSRL & um MDP inteiro do posterior sobre $(P,\rec)$, uma vez por episódio & Osband et al.\ (2013) \\
RLSVI & pesos ruidosos da função valor linear & Osband et al.\ (2016) \\
\emph{Bootstrapped} DQN & uma das $K$ cabeças da rede, uma vez por episódio & Osband et al.\ (2016) \\
\bottomrule
\end{tabular}
\end{center}

\medskip
\noindent
Manter a \emph{mesma} amostra durante todo o episódio é o que produz
\textbf{exploração profunda}: uma sequência coerente de decisões exploratórias,
capaz de alcançar estados que só compensam depois de muitos passos. O
$\epsilon$-guloso, que sorteia de novo a cada passo, executa um passeio
aleatório e nunca chega lá.

\begin{discbox}
Um portal de notícias recebe milhares de acessos por segundo, e com frequência
um novo leitor chega antes de sabermos se o anterior clicou. Nesse regime, o que
se pode dizer da comparação entre amostragem de Thompson e algoritmos de
otimismo? E o que muda se o prior inicial for muito enganoso?

\smallskip
{\footnotesize\color{rlCinza}\emph{Direção da resposta:} explorar a natureza
determinística do UCB (histórico congelado $\Rightarrow$ índice congelado)
contra a natureza estocástica do TS, e separar a pergunta empírica — o que
funciona melhor aqui — da teórica — quem tem a cota mais forte neste cenário.}
\end{discbox}

\begin{discbox}
É legítimo usar amostragem de Thompson para alocar pacientes em um ensaio
clínico? O algoritmo reduz o número de pacientes expostos ao tratamento
inferior, mas produz tamanhos de amostra desbalanceados e dependentes dos dados.

\smallskip
{\footnotesize\color{rlCinza}\emph{Direção da resposta:} contrapor benefício
individual (menos pacientes no braço ruim) a poder estatístico e validade da
inferência posterior; mencionar desenhos adaptativos com restrição de alocação
mínima por braço, e o fato de que a aleatoriedade do TS preserva alguma
proteção contra vieses temporais que um desenho determinístico perderia.}
\end{discbox}

%% ============================================================
\section{Síntese}

\begin{center}\small
\begin{tabular}{@{}lcccc@{}}
\toprule
 & $\epsilon$\textbf{-guloso} & \textbf{UCB} & \textbf{Thompson} & \textbf{Gittins} \\
\midrule
Usa conhecimento prévio        & não & não & sim & sim \\
Determinístico                 & não & sim & não & sim \\
Arrependimento                 & linear se $\epsilon$ fixo & $O(\ln T)$ & $O(\ln T)$ & --- \\
Ótimo dado o prior             & não & não & não & sim$^{\dagger}$ \\
Custo por passo                & $O(m)$ & $O(m)$ & $O(m)$ & alto \\
Estende ao caso contextual     & sim & sim & sim, bem & não \\
Hiperparâmetro de exploração   & $\epsilon$ & constante do bônus & nenhum & --- \\
\bottomrule
\end{tabular}
\end{center}

\noindent{\footnotesize $^{\dagger}$ Sob independência entre braços e critério
descontado em horizonte infinito.}

\medskip
\noindent
Uma leitura útil: o $\epsilon$-guloso explora \emph{cegamente}, o UCB explora
\emph{sistematicamente}, e Thompson explora \emph{na proporção da dúvida}. Os
três exploram; apenas o último deixa que os dados decidam quanto.

\subsection*{Arrependimento e PAC não são a mesma coisa}

Com $\varepsilon = 0{,}05$ no exemplo dos três tratamentos, e definindo o
critério PAC como $\Ind\big(Q(a_t) \ge Q(a^{*}) - \varepsilon\big)$:

\begin{center}\small
\begin{tabular}{@{}ccccccc@{}}
\toprule
\textbf{Otimismo} & \textbf{TS} & \textbf{Ótimo} & \textbf{Arrep.\ otim.} &
\textbf{Otim.\ em $\varepsilon$?} & \textbf{Arrep.\ TS} & \textbf{TS em $\varepsilon$?} \\
\midrule
$a_1$ & $a_3$ & $a_1$ & $0$        & sim & $0{,}85$ & não \\
$a_2$ & $a_1$ & $a_1$ & $0{,}05$   & sim & $0$      & sim \\
$a_3$ & $a_1$ & $a_1$ & $0{,}85$   & não & $0$      & sim \\
$a_1$ & $a_1$ & $a_1$ & $0$        & sim & $0$      & sim \\
$a_2$ & $a_1$ & $a_1$ & $0{,}05$   & sim & $0$      & sim \\
\midrule
\multicolumn{3}{@{}l}{\textbf{Total}} & $\mathbf{0{,}95}$ & $4/5$ & $\mathbf{0{,}85}$ & $4/5$ \\
\bottomrule
\end{tabular}
\end{center}

\medskip
\noindent
Os dois critérios discordam de propósito. O arrependimento penaliza $a_3$
dezessete vezes mais do que $a_2$, porque mede \emph{quanto} se perdeu; o PAC
trata $a_2$ como acerto e $a_3$ como erro, porque só conta \emph{quantas vezes}
se errou feio. Escolher entre eles é decisão de projeto: em um ensaio clínico
importa o dano acumulado; em uma busca de hiperparâmetros, importa terminar com
uma boa configuração.

\subsection*{Formulário}

\begin{center}\small
\begin{tabular}{@{}ll@{}}
\toprule
Arrependimento total & $L_T = \sum_a \E[N_T(a)]\,\Delta_a$, \ \ $\Delta_a = \Vstar - Q(a)$ \\
Regra de Bayes & $p(\theta\mid r) \propto p(r\mid\theta)\,p(\theta)$ \\
Conjugação Beta--Bernoulli & $\Bdist(\alpha,\beta) \to \Bdist(\alpha+S,\ \beta+F)$ \\
Média posterior (Laplace) & $\E[\theta\mid S,F] = (S+1)/(S+F+2)$ a partir de $\Bdist(1,1)$ \\
Casamento de probabilidade & $\pol(a\mid \hist) = \Prob[Q(a) > Q(a')\ \forall a' \neq a \mid \hist]$ \\
Passo de Thompson & $\tilde\theta_a \sim \Bdist(\alpha_a,\beta_a)$;\ \ $a_t = \argmax_a \tilde\theta_a$ \\
Arrependimento bayesiano & $\arrepb(\alg,T) = \E_{\theta\sim p_\theta}[\arrep(\alg,T;\theta)]$ \\
Cota por otimismo & $\arrep \le \E\big[\sum_t U_t(a_t) - Q(a_t)\big]$ \\
\bottomrule
\end{tabular}
\end{center}

%% ============================================================
\section{Exercícios}

\begin{exbox}{11.1 — conjugação para uma sequência}
Mostre, por indução em $n$, que se $\theta \sim \Bdist(\alpha,\beta)$ e
observamos $r_1,\ldots,r_n$ i.i.d.\ $\Bern(\theta)$ com $S = \sum_i r_i$
sucessos e $F = n - S$ fracassos, então
$\theta \mid r_{1:n} \sim \Bdist(\alpha+S,\ \beta+F)$. Conclua que a ordem das
observações é irrelevante e explique por quê.
\end{exbox}

\begin{exbox}{11.2 — casamento de probabilidade em forma fechada}
Suponha $\theta_1 \sim \Bdist(2,1)$, $\theta_2 \sim \Bdist(1,1)$ e
$\theta_3 \sim \Bdist(1,2)$, independentes. Calcule analiticamente
$\Prob[\theta_1 > \theta_2 \text{ e } \theta_1 > \theta_3]$ e as duas
probabilidades análogas para os braços $2$ e $3$. Verifique que somam $1$ e
compare com a frequência com que a amostragem de Thompson escolheria cada braço.
\end{exbox}

\begin{exbox}{11.3 — o guloso tem arrependimento linear}
No exemplo dos três tratamentos, argumente que existe probabilidade
estritamente positiva $p > 0$ de o algoritmo guloso ficar preso em $a_2$ para
sempre, e conclua que $L_T \ge p\,\Delta_2\,(T - c)$ para alguma constante $c$.
Por que o arrependimento médio observado ($13{,}2$ em $T=1000$) parecia
inofensivo?
\end{exbox}

\begin{exbox}{11.4 — cota bayesiana a partir da frequentista}
Prove que $\arrepb(\alg,T) \le \sup_\theta \arrep(\alg,T;\theta)$ e dê um
exemplo (pode ser esquemático) de algoritmo com arrependimento bayesiano
pequeno e arrependimento frequentista grande para algum $\theta$.
\end{exbox}

\begin{exbox}{11.5 — quanto custa um prior errado}
Um braço tem $\theta = 0{,}1$ e prior $\Bdist(100,1)$. Estime quantas puxadas
são necessárias para que a média posterior caia abaixo de $0{,}5$, supondo que
as recompensas saiam na proporção esperada. Compare com o valor observado na
simulação ($204$ puxadas em $T=1000$) e explique a diferença.
\end{exbox}

\begin{exbox}{11.6 — Thompson gaussiano}
Deduza a atualização conjugada Normal--Normal: se
$\theta \sim \Nrm(\mu_0,\sigma_0^2)$ e $r \mid \theta \sim \Nrm(\theta,\sigma^2)$
com $\sigma$ conhecido, mostre que o posterior é gaussiano com precisão
$1/\sigma_1^2 = 1/\sigma_0^2 + 1/\sigma^2$ e média igual à média das duas
médias ponderada pelas precisões. Escreva o laço de Thompson correspondente.
\end{exbox}

\begin{exbox}{11.7 — implementação}
Reproduza a tabela de arrependimento acumulado deste handout comparando guloso,
$\epsilon$-guloso, UCB e Thompson em $T = 2000$, com pelo menos $500$ execuções.
Estenda para $T = 2\times 10^4$ e verifique empiricamente qual curva é
logarítmica e qual é linear. Reporte também a \emph{dispersão} entre execuções,
não só a média.
\end{exbox}

\subsection*{Respostas comentadas}

\textbf{11.1.} Base $n=1$: é a proposição do texto. Passo: pelo posterior ser um
prior para a observação seguinte, $\Bdist(\alpha+S_k,\beta+F_k)$ atualizado por
$r_{k+1}$ dá $\Bdist(\alpha+S_{k+1},\beta+F_{k+1})$. A irrelevância da ordem
decorre de o posterior depender dos dados apenas por $(S,F)$, que é uma
\emph{estatística suficiente} — coerente com as recompensas serem i.i.d.\ dado
$\theta$.

\medskip
\textbf{11.2.} Com densidades $f_1(t) = 2t$, $f_2(t) = 1$, $f_3(t) = 2(1-t)$ e
acumuladas $F_1(t)=t^2$, $F_2(t)=t$, $F_3(t)=2t-t^2$:
\[
  \Prob[\theta_1 \text{ máx}] = \int_0^1 2t\,\cdot t \cdot (2t-t^2)\,dt
  = \int_0^1 (4t^3 - 2t^4)\,dt = 1 - \tfrac{2}{5} = \tfrac{3}{5} = 0{,}60 .
\]
Analogamente,
$\Prob[\theta_2 \text{ máx}] = \int_0^1 t^2(2t-t^2)\,dt = \tfrac{3}{10}$ e
$\Prob[\theta_3 \text{ máx}] = \int_0^1 2(1-t)\,t^3\,dt = \tfrac{1}{10}$.
Somam $1$, como devem. Uma simulação com $4\times10^5$ amostras dá $0{,}598$,
$0{,}301$ e $0{,}101$ — e essas são \emph{exatamente} as frequências com que
Thompson escolheria cada braço, pela proposição da Seção~4.

\medskip
\textbf{11.3.} Basta que, nas primeiras puxadas, $a_1$ falhe e $a_2$ acerte:
esse evento tem probabilidade positiva ($\approx 0{,}05 \times 0{,}90$ na
inicialização usada), e a partir daí a média empírica de $a_2$ nunca é superada,
porque $a_1$ deixa de ser puxado. Nesse evento o arrependimento cresce a taxa
$\Delta_2$ por passo, o que dá a cota. O arrependimento médio parecia inofensivo
porque $\Delta_2 = 0{,}05$ é minúsculo: a reta tem inclinação pequena. Ao
quadruplicar $T$, o arrependimento do guloso quadruplica ($58{,}9 \to 229{,}9$),
enquanto o de Thompson cresce como $\ln T$ ($13{,}0 \to 15{,}7$).

\medskip
\textbf{11.4.} A primeira parte é a monotonicidade da esperança. Para a segunda,
considere um algoritmo que ignore os dados e jogue sempre o braço com maior
valor \emph{a priori}: se o prior concentra quase toda a massa num $\theta$ em
que esse braço é ótimo, o arrependimento bayesiano é quase nulo, mas para o
$\theta$ (de probabilidade a priori pequena) em que o braço é péssimo o
arrependimento frequentista é linear em $T$.

\medskip
\textbf{11.5.} A média posterior é $100/(100+1+F)$, com $F$ fracassos
acumulados; ela cai abaixo de $0{,}5$ quando $F > 99$. Como $\theta = 0{,}1$,
cerca de $90\%$ das puxadas são fracassos, e são necessárias
$\approx 99/0{,}9 \approx 110$ puxadas. A simulação registra $204$ porque o
critério real não é a média posterior cair abaixo de $0{,}5$, e sim as
\emph{amostras} do posterior deixarem de ganhar do outro braço — o que exige
concentração adicional, já que a Beta ainda tem cauda à direita.

\medskip
\textbf{11.6.} Multiplicando as duas densidades gaussianas e completando o
quadrado em $\theta$, os termos quadráticos dão
$\sigma_1^{-2} = \sigma_0^{-2} + \sigma^{-2}$ e os lineares dão
$\mu_1 = \sigma_1^2(\mu_0/\sigma_0^2 + r/\sigma^2)$. O laço de Thompson passa a
sortear $\tilde\theta_a \sim \Nrm(\mu_a, \sigma_a^2)$ e a atualizar
$(\mu_{a_t}, \sigma_{a_t})$ pela fórmula acima. Note a interpretação: precisões
se somam, e a média nova é a média ponderada pelas precisões — o dado ``puxa''
o prior tanto mais quanto menos ruidoso for.

\medskip
\textbf{11.7.} Exercício de implementação; o esqueleto de seis linhas da
Seção~4 basta para o TS. Um resultado esperado: a dispersão do guloso é
bimodal (execuções que acertam o braço bom têm arrependimento quase nulo;
as que travam têm arrependimento linear), enquanto a do TS é unimodal e
estreita. Essa diferença de \emph{forma} da distribuição é invisível na média
e é o argumento mais forte contra o guloso.

%% ============================================================
\section{Autoavaliação e glossário}

\noindent Você domina esta aula se consegue, sem consultar o material:

\begin{multicols}{2}
{\footnotesize
\begin{itemize}[leftmargin=1.5em, itemsep=0.3ex, topsep=0.2ex, label=$\square$]
  \item escrever $L_T = \sum_a \E[N_T(a)]\Delta_a$ e explicar por que ela
        organiza toda a análise;
  \item enunciar e demonstrar a conjugação Beta--Bernoulli;
  \item ler $(\alpha,\beta)$ como pseudo-contagens e recuperar a regra de
        Laplace;
  \item definir casamento de probabilidade e dizer por que é difícil de
        calcular;
  \item escrever o laço de Thompson e implementá-lo em poucas linhas;
  \item demonstrar que Thompson implementa casamento de probabilidade;
  \item distinguir arrependimento frequentista de bayesiano e provar a
        desigualdade entre eles;
  \item explicar o papel do otimismo como \emph{técnica de demonstração};
  \item dar um exemplo concreto em que um prior enganoso arruína o TS;
  \item explicar por que o guloso tem arrependimento linear mesmo parecendo
        bom em horizonte curto;
  \item enunciar o teorema do índice de Gittins e suas três limitações;
  \item explicar o que é exploração profunda e por que $\epsilon$-guloso não a
        produz.
\end{itemize}}
\end{multicols}

\subsection*{Convenções de tradução fixadas nesta aula}

\begin{center}\small
\begin{tabular}{@{}ll@{}}
\toprule
\textbf{Inglês} & \textbf{Português} \\
\midrule
regret / Bayesian regret & arrependimento / arrependimento bayesiano \\
arm / gap & braço / lacuna \\
prior / posterior & prior / posterior \\
conjugate prior & prior conjugado \\
probability matching & casamento de probabilidade \\
Thompson sampling / posterior sampling & amostragem de Thompson / amostragem posterior \\
optimism in the face of uncertainty & otimismo diante da incerteza \\
index policy / Gittins index & política de índice / índice de Gittins \\
contextual bandit & bandit contextual \\
belief-state MDP & MDP de crenças \\
deep exploration & exploração profunda \\
\bottomrule
\end{tabular}
\end{center}

\medskip
\noindent{\footnotesize\color{rlCinza}
\textbf{Leituras.} Sutton e Barto, \emph{Reinforcement Learning: An
Introduction}, 2.\textsuperscript{a} ed., cap.~2 (seções 2.7 e 2.9)
\textbullet\ Russo, Van Roy, Kazerouni, Osband e Wen, \emph{A Tutorial on
Thompson Sampling} (2018) \textbullet\ Lattimore e Szepesvári, \emph{Bandit
Algorithms} (2020), caps.~34--36 \textbullet\ Chapelle e Li, \emph{An Empirical
Evaluation of Thompson Sampling} (NeurIPS 2011) \textbullet\ Russo e Van Roy,
\emph{Learning to Optimize via Posterior Sampling} (2014).

\medskip
Material adaptado e expandido de \textbf{CS234: Reinforcement Learning}, Emma
Brunskill, Stanford University. Prof.\ Marcelo Keese Albertini, Faculdade de
Computação, Universidade Federal de Uberlândia.}

\end{document}
